% !TEX program = lualatex
\documentclass[aspectratio=43,professionalfonts,handout]{beamer}
\usepackage{beamer_template}

\newcommand{\LectureNum}{19}
\newcommand{\LectureTitle}{グラフの拡大と縮小（１）}
\newcommand{\TextPages}{p.\,152-154}
\newcommand{\NextTopic}{グラフの拡大と縮小（２）}
\newcommand{\NextPages}{p.\,154}
\newcommand{\CourseName}{基礎数学B}
\renewcommand{\TermName}{後期}

\begin{document}

% -----------------------------------------------
% スライド1: 表紙＋目標＋用語・定義
% -----------------------------------------------
\begin{frame}{\CourseName\quad 第\LectureNum 回「\LectureTitle 」}
  {\small \TermName\quad \TeacherName\quad ／\quad 教科書 \TextPages\quad
  \textbf{目標}：$y = Cf(x)$, $y = f(Cx)$ 型のグラフ変換規則を理解できる}

  \begin{block}{用語・定義}
  \begin{description}[labelwidth=5.5em]
    \item[\kw{振幅：$|a|$}] $y = a\sin x$ の振れ幅は $|a|$（$y$ 方向のスケール）
    \item[\kw{周期：$2\pi/|b|$}] $y = \sin(bx)$ の周期は $\dfrac{2\pi}{|b|}$（$|b| > 1$ で縮小，$0 < |b| < 1$ で拡大）
    \item[\kw{$b > 1$ で縮小、$0 < b < 1$ で拡大}] グラフが $x$ 方向に縮む（または伸びる）
  \end{description}
  \end{block}
\end{frame}

% -----------------------------------------------
% スライド2: 公式・性質
% -----------------------------------------------
\begin{frame}{公式・性質}
  \begin{block}{}
  \begin{itemize}
    \item $y = a\sin x$：振幅 $|a|$、周期 $2\pi$
    \item $y = \sin(bx)$：振幅1、周期 $2\pi/b$
    \item $y = a\sin(bx)$：振幅 $|a|$、周期 $2\pi/b$
  \end{itemize}
  \end{block}
\end{frame}

% -----------------------------------------------
% スライド3: 例1→問1
% -----------------------------------------------
\begin{frame}{例1→問1}
  \begin{exampleblock}{例1) p.162（振幅・周期を答えグラフをかけ）}
    （板書で解説）
  \end{exampleblock}

  \vfill

  \begin{block}{問1) p.162\quad 自分で解いてみよう}
    （ノートに解くこと）
  \end{block}
\end{frame}

% -----------------------------------------------
% スライド4: 例2→問2
% -----------------------------------------------
\begin{frame}{例2→問2}
  \begin{exampleblock}{例2) p.163（$y = \cos(2x)$ のグラフ）}
    （板書で解説）
  \end{exampleblock}

  \vfill

  \begin{block}{問2) p.163\quad 自分で解いてみよう}
    （ノートに解くこと）
  \end{block}
\end{frame}

% -----------------------------------------------
% まとめ
% -----------------------------------------------
\begin{frame}{まとめ}
  \begin{itemize}
    \item $y = a\sin(bx)$：振幅 $|a|$，周期 $\dfrac{2\pi}{|b|}$
    \item $y = a\cos(bx)$ も同様の規則
    \item $b$ が大きいほど波が細かくなる（周期が短い）
  \end{itemize}

  \vfill
  {\small 基本問題：なし（演習問題集で確認）}
\end{frame}

\end{document}
